\scanpage{53}
% archive sheet number: 52
Supposons maintenant $k$ isomorphe à $\mathbb C$, et \unsure{bornons|donnons} nous aux places pour lesquelles $k$ est \uline{complet}. Alors on a des réseaux
\[
M_v^{\mathbb Z}\subset M^{\mathbb A_{\mathbb C}}\qquad\text{où}\quad M^{\mathbb A_{\mathbb C}}=\deleted{\unsure{\mathbb A_{\mathbb C}\otimes_{\mathbb Z}M}}\prod_{l\ \mathrm{fini}\ \inserted{\text{ou }\uline{\text{infini}}}} M(\ell)
\]
est indépendant de $v$.

\deleted{On voit encore}

Si $M$ est défini sur \inserted{le corps de type fini} $k_0$, et $k_0$ assez \unsure{petit|gros} pour que $G_M$ connexe, mais choisi de degré de transcendance minimum, \deleted{alors} \deleted{si $\sigma\in\operatorname{Aut}(k/k_0)$ \unsure{tel que} $\sigma$} \deleted{\unsure{du}} \deleted{\unsure{pt} \unsure{est}} si $\sigma\in\deleted{\operatorname{Aut}_{k_0}k}\operatorname{Gal}(k/k_0)$
\margin{\uline{NB} $k$ \uline{n'est pas} \deleted{\unsure{algébrique}} \uline{alg. sur $k_0$} mais de deg. de transc. infini}
on a
\[
M_{\sigma v}^{\mathbb Z}=\dot{\sigma}M_v^{\mathbb Z}
\]
où $\dot{\sigma}$ opère sur $M^{\mathbb A_{\mathbb C}}$ \deleted{\unsure{en utilisant}} \deleted{\unsure{son image|les images}} de la façon suivante
\[
\dot{\sigma}\bigl((x_l)\unsure{\}}_{l\ \mathrm{fini}},x_\infty\bigr)=(\sigma_l x_l,\dot{\sigma}_\infty x_\infty)
\]
où $\sigma_l$ provient de l'hom canonique $\operatorname{Gal}(k/k_0)\unsure{\mathrel{\text{---}}|\longrightarrow}G(\mathbb Z_l)$, et $\dot{\sigma}_\infty$ est l'hom \uline{$\frac12$ linéaire} défini par $\sigma$ via
\[
M_{\unsure{/k|k}}(\infty)\simeq M_{k_0}(\infty)\otimes_{k_0} k
\]
